\section{Algorithm}
The central algorithm fits a penalized trend for each candidate Guerrero
smoothness setting, evaluates out-of-sample validation forecasts, and then
refits the selected model before test evaluation. This section summarizes the
mechanism behind the penalized estimator, the numerical linear algebra used to
fit it, and the validation search used to choose among candidates. The
penalized objective belongs to the same finite-difference smoothing family as
Whittaker-Henderson graduation, Guerrero smoothing, and penalized spline
methods \citep{whittaker1922graduation,henderson1924graduation,
guerrero2007penalized,eilers1996psplines,eilers2003perfect,wahba1990splines}.

Let \(z=(z_1,\ldots,z_N)^\top\) denote the observed training series, let \(d\)
be the finite-difference order, and let
\[
  K \in \mathbb{R}^{(N-d)\times N}
\]
be the order-\(d\) differencing matrix. For a penalty value \(\lambda > 0\),
the penalized trend estimate is defined by
\begin{equation}
\label{eq:algorithm_pls}
  \hat{t}(\lambda)
  =
  \arg\min_{t\in\mathbb{R}^N}
  \left\{
    \|z-t\|_2^2 + \lambda \|K t\|_2^2
  \right\}.
\end{equation}
The first term measures fidelity to the observed series and the second term
penalizes roughness. Larger values of \(\lambda\) produce smoother fitted
trends.

The objective in \eqref{eq:algorithm_pls} is quadratic. Define
\[
  A(\lambda) = I_N + \lambda K^\top K.
\]
Taking the derivative with respect to \(t\) and setting it equal to zero gives
\[
  A(\lambda)\hat{t}(\lambda)=z.
\]
Since \(A(\lambda)\) is positive definite for \(\lambda>0\), the fitted trend is
\begin{equation}
\label{eq:algorithm_smoother}
  \hat{t}(\lambda)
  =
  S(\lambda)z,
  \qquad
  S(\lambda)=\left(I_N+\lambda K^\top K\right)^{-1}.
\end{equation}
Here \(S(\lambda)\) is the smoothing matrix.

\subsection{Spectral solution of the penalized system}
Although \eqref{eq:algorithm_smoother} is written with an inverse, the
implementation does not form a dense inverse directly. For a fixed sample size
\(N\) and order \(d\), it first constructs the penalty matrix
\[
  P_d = K^\top K.
\]
This matrix is symmetric positive semidefinite, so it admits an orthogonal
eigendecomposition
\[
  P_d = Q \Delta Q^\top,
  \qquad
  \Delta=\mathrm{diag}(\delta_1,\ldots,\delta_N).
\]
Then
\[
  I_N+\lambda P_d
  =
  Q(I_N+\lambda \Delta)Q^\top,
\]
and the fitted trend can be computed as
\begin{equation}
\label{eq:algorithm_spectral_solution}
  \hat{t}(\lambda)
  =
  Q
  \mathrm{diag}\left(
    \frac{1}{1+\lambda \delta_1},
    \ldots,
    \frac{1}{1+\lambda \delta_N}
  \right)
  Q^\top z.
\end{equation}
This reduces each candidate fit to projections onto the eigenvectors, elementwise
division by \(1+\lambda\delta_i\), and a projection back to the original time
domain. Since the same \(Q\) and \(\delta_i\) are reused for many candidate
\(\lambda\) values at the same \(N\) and \(d\), the paper pipeline benefits from
caching the solver for repeated fits. The same eigenvalues also provide the
effective-degrees-of-freedom trace used in Guerrero's index, so the numerical
linear algebra for fitting and the numerical linear algebra for tuning are
shared rather than duplicated.

For Guerrero-style penalized trends, the implementation may also estimate a
constant \(d\)-th difference drift. In that case the right-hand side is adjusted
iteratively and the drift estimate \(m\) is updated from the fitted trend until
the change is smaller than a numerical tolerance. HPTrend and WhittakerTrend are
run as classical fixed-drift baselines with this drift set to zero.

\subsection{Mapping Guerrero smoothness to \texorpdfstring{\(\lambda\)}{lambda}}
The paper reports Guerrero smoothness for finite-difference penalized smoothers,
but the linear system is solved in terms of \(\lambda\). The eigendecomposition
above gives a numerically convenient expression for the trace of the smoothing
matrix:
\[
  \mathrm{tr}\{S(\lambda)\}
  =
  \sum_{i=1}^N \frac{1}{1+\lambda\delta_i}.
\]
For \(d\ge 1\), the implementation uses the normalized Guerrero index
\begin{equation}
\label{eq:algorithm_guerrero_normalized}
  s(\lambda;N,d)
  =
  \frac{
    1 - N^{-1}\sum_{i=1}^N (1+\lambda\delta_i)^{-1}
  }{
    1-d/N
  },
\end{equation}
which maps the attainable range to a scale close to \([0,1)\) for the finite
sample and difference order \citep{guerrero2007penalized}. Given a requested
Guerrero smoothness value \(s\), the code finds the corresponding \(\lambda\)
by bisection:
\begin{enumerate}
  \item Set an initial bracket \([0,1]\).
  \item Increase the upper endpoint by factors of ten until the target
        smoothness is bracketed.
  \item Bisect the interval until the trace-based smoothness is within the
        numerical tolerance.
\end{enumerate}
This step is monotone because increasing \(\lambda\) decreases the effective
degrees of freedom \(\mathrm{tr}\{S(\lambda)\}\) and therefore increases the
Guerrero smoothness index.

\subsection{How Guerrero's idea defines the penalized methods}
Guerrero's idea is not a separate estimator in the pipeline; it is the
trace-based scale used to parameterize and compare finite-difference penalized
smoothers. Once \(N\) and \(d\) are fixed, every penalty value \(\lambda\)
corresponds to a Guerrero smoothness value through the smoothing matrix trace.
The implementation uses this correspondence in four related ways.

For TrainValidationSelector, the candidate set is written directly as pairs
\((d,s)\), where \(s\) is the Guerrero smoothness index. Each candidate \(s\)
is converted to its corresponding \(\lambda\), the penalized system is solved
on the training segment, and the resulting trend is extrapolated into the
validation segment. The selected model is the candidate with the smallest
validation error. In this sense, Guerrero's index becomes the hyperparameter
that the temporal validation split chooses.

For TimeWeightedValidationSelector, the same Guerrero-indexed candidate set is
used. The difference is only in the validation criterion: later validation
observations receive larger weights, so the selected Guerrero smoothness is the
one that performs best for the more recent part of the validation window. The
underlying finite-difference penalty, the \(s\)-to-\(\lambda\) conversion, and
the linear solve are otherwise the same as in TrainValidationSelector.

For HPTrend, the trend is also a finite-difference penalized least-squares
smoother, usually with a second-difference penalty. The classical HP
presentation specifies a penalty \(\lambda\) rather than a Guerrero smoothness
index. In the paper pipeline, that fixed \(\lambda\) is passed through the same
trace formula to report the implied Guerrero smoothness. HPTrend is therefore
Guerrero-compatible for interpretation, even though the validation set is not
used to choose its penalty.

For WhittakerTrend, the fitted trend again solves the same algebraic family of
linear systems, with a finite-difference penalty and a fixed penalty setting.
As with HPTrend, the fixed \(\lambda\) determines a fitted smoother matrix and
therefore an implied Guerrero smoothness index. This lets WhittakerTrend be
reported beside the validation-selected methods without pretending that it ran
a validation search.

This is why the paper reports Guerrero smoothness for TrainValidationSelector,
TimeWeightedValidationSelector, HPTrend, and WhittakerTrend, but leaves that
field blank for moving averages, exponential smoothing, and polynomial
forecasting. Those latter methods can still be evaluated by forecast error and
realized roughness, but they do not have a native \(S_d(\lambda;N)\) because
they are not defined by the penalized system in
\eqref{eq:algorithm_pls}.

For validation, let \(R\) be the matrix that selects validation observations and
let \(n_{\mathrm{val}}\) be the number of validation points. The validation
residual is
\[
  r(\lambda) = Rz - R S(\lambda)z,
\]
and the validation mean squared error is
\begin{equation}
\label{eq:algorithm_validation_mse}
  f(\lambda)
  =
  \frac{1}{n_{\mathrm{val}}}
  r(\lambda)^\top r(\lambda).
\end{equation}
The derivative calculation in the companion report follows from
\[
  S'(\lambda)
  =
  -S(\lambda)K^\top K S(\lambda),
\]
which implies
\[
  r'(\lambda)
  =
  R S(\lambda)K^\top K S(\lambda)z.
\]
Therefore,
\begin{equation}
\label{eq:algorithm_validation_derivative}
  f'(\lambda)
  =
  \frac{2}{n_{\mathrm{val}}}
  \left[ R\left(I_N-S(\lambda)\right)z \right]^\top
  \left[ R S(\lambda)K^\top K S(\lambda)z \right].
\end{equation}

\subsection{Validation search and local minima}
The derivative in \eqref{eq:algorithm_validation_derivative} explains the
shape of the validation objective, but the implemented paper pipeline uses a
robust direct search over candidate Guerrero smoothness values. For each
candidate difference order \(d\) and smoothness value \(s\), the algorithm:
\begin{enumerate}
  \item converts \(s\) to \(\lambda\) using the bisection procedure above;
  \item fits the penalized trend on the training segment only;
  \item extrapolates the fitted trend into the validation segment;
  \item computes the validation loss between the forecast and the observed
        validation log prices.
\end{enumerate}
For the S\&P 500 comparison, the paper script evaluates orders
\(d\in\{1,2,3\}\) and a fixed grid of Guerrero smoothness values between 0.05
and 0.98. The broader library default uses a denser grid, but the paper uses a
smaller reproducible grid to keep the empirical study fast and stable.

The selector records the full validation curve and detects local minima on the
one-dimensional smoothness grid for each order. A grid point is treated as a
local minimum when its validation score is no larger than the scores of its
neighbors; boundary points can also be local minima if they improve on their
single neighbor. The library can optionally refine each detected bracket by
golden-section search. In the paper pipeline, refinement is disabled so that the
saved validation curves, local minima, and selected value all correspond exactly
to the stated reproducible grid.

The time-weighted selector uses the same candidate search, but replaces the
ordinary mean squared validation loss with a weighted validation loss. In the
paper run, the weights follow an exponential decay scheme, so later validation
observations receive relatively more emphasis than earlier validation
observations. This changes only the scoring rule, not the penalized smoothing
linear algebra. In forecasting terminology, this is a validation-rule change
rather than a different smoothing model \citep{makridakis1998forecasting,
hyndman2008exponential}.

\subsection{Forecasting from a fitted penalized trend}
After fitting a penalized trend on the training segment, the package forecasts
future trend values by extending the finite-difference structure at the right
tail. If the selected order is \(d\), the forecast recursively imposes a
constant estimated \(d\)-th difference \(m\). For \(d=1\), this is a constant
slope extension; for \(d=2\), it is a constant second-difference extension; and
for higher orders it extends the corresponding finite-difference recurrence.
These forecasts are used for validation scoring and, after refitting on
training plus validation, for final test evaluation. The benchmark forecasts
are intentionally simpler: trailing moving averages, exponential smoothing, and
polynomial extrapolation are used as transparent baselines rather than as a
complete forecasting taxonomy \citep{brown1959forecasting,holt2004exponential,
winters1960forecasting,gardner1985exponential,gardner2006exponential,
draper1998regression}.

\subsection{Final comparison protocol}
The complete numerical procedure for selector-based penalized methods is:
\begin{enumerate}
  \item split the series temporally into training, validation, and test
        segments;
  \item fit and score each \((d,s)\) candidate on training-to-validation
        forecasts;
  \item choose the candidate with the lowest validation loss, without using the
        test segment;
  \item refit the chosen penalized smoother on training plus validation;
  \item forecast the test segment and compute MAE, RMSE, SMAPE, and related
        metrics;
  \item compute realized roughness \(R_d(\hat{t})=\|D^d\hat{t}\|_2^2\) from the
        fitted training-plus-validation trend.
\end{enumerate}
The non-Guerrero baselines use the same temporal protocol, but they do not have
a Guerrero smoothness search. The moving-average baseline uses its fixed
trailing window, the exponential-smoothing baseline uses its fixed \(\alpha\),
and the polynomial forecaster uses its fixed degree. They still receive test
forecasts and realized roughness diagnostics, but their Guerrero smoothness
entries are left blank.

The Guerrero smoothness index used by the implementation is
\[
  S_d(\lambda; N)
  =
  1 -
  \frac{\mathrm{tr}\left[S(\lambda)\right]}{N},
  \qquad
  S(\lambda)=\left(I_N+\lambda K^\top K\right)^{-1}.
\]
This index is meaningful for the finite-difference penalized smoothers in the
comparison: the train-validation selector, the time-weighted selector, HPTrend,
and WhittakerTrend. It is not a native parameter of moving averages,
exponential smoothing, or polynomial forecasting.

In practice, solving \(f'(\lambda)=0\) does not give a simple general closed
form. The implementation therefore uses numerical validation search over the
package's normalized Guerrero smoothness parameter, converts the selected index
to \(\lambda\), and then refits the final penalized trend on the combined
training and validation sample. The comparative pipeline applies the same
temporal selection principle to the proposed train-validation selector and its
time-weighted variant, while fixed-parameter baselines are fit on the same
training and refit samples.
